% Proposed complete written endpoint assembly, 11 October 2026.
% This source proof is under review. It is not a new Lean endpoint theorem,
% a clean replay, or an independent-checker certificate.
\section{The endpoint and the target-uniform continuation argument}
\label{sec:endpoint-assembly}

We now insert the local physical moment into the pinned numerator
construction. The input results are the established incumbent moments,
the exact physical Euler expansion and coefficient mass in B0.32--B0.44
and B4 of the pinned Cycle25 manuscript, its detector count B0.40, and
its contour and continuation construction B0.7 and T2--T3. The new input
is the local-moment proposition proved in
\path{PhysicalLocalAssembly.tex}, with the fixed-dictionary and
height-uniform density arguments preceding it. The downward-scale
argument is \path{PhysicalScaleShortening.tex}. We retain the exact
rational certificate in \path{local_moment_endpoint_audit.py}.
No unproved varying-parent moment is introduced in this application.

Set
\[
 \theta=\frac{4374433}{5000000},\quad
 \ell=\frac{626701}{3750000},\quad
 h=\frac{3969737987}{5000000000},\quad
 x=1+\ell-h,\quad y=h-2\ell,
 \qquad C=\frac{2249}{500}.
\]
In particular $\theta=11/12-\ell/4$, $x+y=1-\ell$,
$0<\ell<x<y<2x$, and $y<1/2<h<2y$.

\begin{theorem}[Proposed endpoint, relative to the pinned analytic framework]
Every finite-order Hecke character of the Eisenstein field with trivial
infinite type has no zero in $\Re s>\theta$, with the principal pole at
$s=1$ excluded. Presented characters may be imprimitive and have
arbitrary conductor. Every ordinary Dirichlet $L$-function has the same
zero-free half-plane, again excluding the principal pole. In particular
$\zeta(s)\ne0$ there for $s\ne1$.
\end{theorem}

\begin{proof}[Proposed source proof, under review]
Let $\beta_*$ be the source's Hecke zero supremum, including its sentinel
$1/2$. The incumbent theorem gives $\beta_*\le7/8$. Suppose
$\mu=\beta_*-\theta>0$ and set
\[
 \kappa=2\beta_*-1,\qquad
 \kappa_*=2\theta-1=\frac{1874433}{2500000}.
\]
Then $\kappa_*<\kappa\le3/4$. The incumbent marked detector and ordinary
moment apply at this $\kappa$: their hypothesis $2\beta_*-1\le\kappa$
holds by equality. The new local moment uses only the incumbent
$9/2$ and unrestricted zero-slot estimates at descendant states.

\paragraph{1. The actual physical family and selected slots.}
Fix a primitive target $\eta$ only after the common real margins and slot
mesh have been chosen. Its finite exclusion set $S$, original local
Euler factors, and all natural row zeros are those of the source probe.
In the physical numerator moment the fixed base is
\texttt{fixedSourcePrincipal $S$}; the variable part is the sextic row
character. This is exactly the coefficient identity in
\path{PhysicalMomentInput.source_physical_slot_conj} and
\path{reflected_pair_norm_eq}. The relative prime coefficient
contains the inverse fixed base and the common conjugate annular
profile, with real power $33/50$ and individual height $-\Im z_i$.
Conjugation applies to the entire reflected product. Thus the local
proposition applies to this fixed family with the same literal
eligibility and masks. The fixed exceptional free rows have bounded
norm and are absent from every growing retained dyad; the source
\path{weighted_source_exceptional_norm} proves this fact.

In an amplitude class let $w_i=\ell_i/d$, let
$L=\sum_iw_i=\ell/d$, and let $0\le g_i\le\delta/2$ be its amplitude
exponents. Error slots have $g_i=0$. Put
$q=L^{-1}\sum_iw_ig_i=\delta p$, where $0\le p\le1/2$.
For any target capacity $0\le z_A\le L$, the exact theorem
\path{SlotSelection.zero_error_whole_slots} selects positive-gain
main slots $I$ such that
\[
 z_I=\sum_{i\in I}w_i\le z_A,
 \qquad
 \left|\prod_{i\in I}Q_i\right|^2
 \ge U^{2qz_A-\delta\,\mathrm{mesh}-o(1)}.       \tag{EA.1}
\]
The gain is measured at the target capacity $z_A$, even if $z_I$ is
strictly smaller. It is not legitimate to assume $z_I$ nearly fills
$z_A$ or to insert error slots into the selected product. The selection
is made once for the amplitude class, not independently for its rows.

Put $d_0=1/200$ and $\zeta_0=10^{-10}$. All original slots can be fixed
so short that their relative widths meet every required mesh for
$d\ge d_0$. Detector supply is available because the exact geometry
satisfies $\ell/(h+\zeta_0)>4/21$.
For $d<1/2$, the same selection theorem applies directly, without the
$d\ge1/2$ assumption of its convenience wrapper:
$L\ge2\ell>1/3$ and the old physical capacity is at most
$1/(6\kappa)\le5/21$. Thus no lower-row supply condition is missing.

\paragraph{2. The reflected scale and every free Euler shift.}
The exact coefficient expansion writes a shift as $h_0h_r$, with
$h_r$ charged to the repeated-prime conductor deficit. B0.42--B0.44 give
the actual reflected scale bound
\[
 X_u\ll_S (U/Y)h_0,
 \qquad \sum_{h_0,h_r}|c_u(h_0,h_r)|h_0^{\delta/2}
       \ll U^\varepsilon.                            \tag{EA.2}
\]
The same combined conductor occurs in both statements. Its deficit is
not charged a second time. Natural Euler restoration only shortens
scales; its coefficients have the source's bounded absolute mass.
The original support retains $h_0\ll Y$, so all scale exponents lie in
a fixed bounded range.

Consider first $h-10^{-3}\le d\le h+\zeta_0$ and
$t=\log_Uh_0\in[0,10^{-3}]$. The upper reflected length is
$n(d,t)=1-y/d+t$. Choose the target capacity just below
\[
 z_A(t)=\frac{1-2n(d,t)}{C},
\]
reserving the small capacity decrement for annular and scale-box errors.
The exact geometric checks, with strict room for those errors, give
\[
 \frac{41}{100}<n(d,t)\le\frac{211}{500},\qquad
 \frac3{100}<z_A(t)\le\frac{71}{2000},\qquad z_A(t)<L.
\]
Let $x_1,x_2\le n(d,t)$ be the actual nonnegative plain lengths, and
let $z_I$ be the actual selected width from (EA.1). There are three cases.
If $z_I<3/100$, the incumbent capacity holds with strict reserve:
\[
 x_1+x_2+\frac92z_I
 \le2\frac{211}{500}+\frac92\frac3{100}
 =\frac{979}{1000}<1.                                \tag{EA.3}
\]
If $z_I\ge3/100$ and some $x_j<41/100$, the downward-scale lemma gives
the incumbent reserve $33/4000$. In the remaining case the local
proposition applies because
$x_1+x_2+Cz_I\le2n(d,t)+Cz_A(t)\le1$.
The empty selected product uses the unrestricted zero-slot estimate.
The same two-variable scale-supremum recovery and its four derivative
profiles apply on these separate regions. Therefore the fourth moment
needed for numerator H\"older has the same $U^{1+\varepsilon}$ bound
in all cases, while its gain remains $qz_A(t)/2$ by (EA.1).

For these small shifts,
\[
 -\frac{\delta t}{2}-\frac{qz_A(t)}2
 \le-\frac{qz_A(0)}2,
\]
since $q\le\delta/2$ and $C>1$. For larger shifts use the incumbent
capacity instead. Write $b=2y/d-1$ and use the smaller comparison
capacity $z_o(t)=\max(0,(b-2t)/(9/2))$. The exact scalar audit shows
\[
 S(10^{-3})>\frac3{10000},\qquad
 S(t)=\frac t2-\frac{b/C-z_o(t)}4.
\]
While $z_o(t)>0$, its derivative is $7/18$; after the capacity clips
to zero the derivative is $1/2$. Hence $S(t)\ge0$ for every
$t\ge10^{-3}$. With $q\le\delta/2$ this proves
\[
 -\frac{\delta t}{2}-\frac{qz_o(t)}2
 \le-\frac{qz_A(0)}2.                                \tag{EA.4}
\]
The actual old capacity at coefficient $6\kappa\le9/2$ is no smaller.
Thus (EA.2) pays every free-shift dyad, including those with zero old
capacity. No Euler term has been omitted from the numerator inversion.
At zero capacity the two reflected plain lengths can have sum greater
than one; this is exactly where the \emph{unrestricted} physical MP0 is
needed. Its actual source adapter
\path{Physical/CertifiedZero.finite_pair_zeroAt} has no affine
plain-length condition and bounds the moment by $U^{1+\varepsilon}$
for each plain scale up to $U^2$. The theorem
\path{Physical/OriginalFourth.certified_original_numerator_fourth}
accordingly accepts either its affine capacity inequality or an empty
selected family. The original support bound in (EA.2) gives reflected
length at most $1+o(1)<2$, so this is a literal application of that
theorem. The weaker all-free-row plain estimate used in the new density
argument is not being substituted for this physical MP0.
Below $d=h-10^{-3}$ the original B0.41 argument is unchanged.

\paragraph{3. The common retained-row saving.}
The detector remains
\[
 R_\kappa(\delta,p)=1-\delta+
 \frac{\delta c_\kappa(p)(5/6-\delta)}
      {2(5/6-\delta+\delta c_\kappa(p))},\qquad
 c_\kappa(p)=\frac{2(1-p)(1-p/(3\kappa))}
                   {3-p-2p/(3\kappa)}.
\]
All $\delta\in[0,3/4]$ and $p\in[0,1/2]$ are included. The physical
H\"older argument and the weighted mass (EA.2), on the original full
$w$-inversion line, give the exponent $G_{C,\kappa}$ of the paper on
the local row interval. On old rows its general form is
\[
 G_o(d)=\frac{1+\delta}{2}+h\frac{13}{75}
 -\frac y2-\frac\ell2+\ell q
 +d\left(\frac14+\frac{3R_\kappa}4-\frac{17}{50}\right)
 -\frac{q}{12\kappa}\min(d,2y-d).                    \tag{EA.5}
\]
This is the exact B0.46/B0.55 ledger, retaining the reciprocal only
after the full physical numerator inversion.

At $\kappa=\kappa_*$ the rational Bernstein certificate proves savings
$10^{-7}$ at $(d,c)=(h,C)$, $10^{-5}$ at
$(h-10^{-3},6\kappa_*)$, and $10^{-3}$ at
$(1/2,6\kappa_*)$, uniformly on the full $(\delta,p)$ rectangle.
At the junction the new exponent is no larger than the old one because
$C<6\kappa_*$ and $2y-d>0$. Affinity in $d$ covers the intervening
local interval. The extension above $h$ costs at most $2\zeta_0$, so
the retained saving there is at least $499/5000000000>9\cdot10^{-8}$.

The old function (EA.5) is convex in $d$. To cover $d_0\le d\le1/2$,
combine the certificate at $1/2$ with the elementary bound at $d_0$,
using $a\le7/8$, $R_\kappa\le1$ and $q\le3/8$:
\[
 G_o(d_0)\le\frac78-\frac{49h}{150}+\frac{7\ell}{8}
                  +\frac{33d_0}{50}
 =\frac{573880513637}{750000000000}<\theta-\frac1{10}.
\]
On $[1/2,h-10^{-3}]$, affinity and its two certified endpoints suffice.
The floor classes use the source bound
\[
 \frac{51}{100}+\frac h3+\frac{51\ell}{100}
 =\frac{12898208027}{15000000000}<\theta-\frac1{100},
\]
with room for the tiny row extension. They do not require a positive
detector zero above the sentinel.

The variation in $\kappa$ must be checked separately for the two
physical capacities. On the local branch only $R_\kappa$ varies, and
the source derivatives $\partial_\kappa c_\kappa\le1/18$ and
$\partial_cR\le5/12$ give $\partial_\kappa G_{C,\kappa}<1/50$.
Because $\kappa-\kappa_*=2\mu$, its increase is less than
$\mu/25$. Thus its saving below $\beta_*$ is at least its certificate
saving plus $24\mu/25$.
For the old branch the reciprocal capacity also varies. Since
$\min(d,2y-d)\le1/2$ and $d\le1$, one has
\[
 \partial_\kappa G_o
 \le\frac5{288}+\frac{(3/8)(1/2)}{12(7/10)^2}
 =\frac{695}{14112}<\frac18.
\]
Its feedback saving is therefore at least its certificate saving plus
$3\mu/4$. In particular all retained classes have a common strictly
positive power reserve before a target character is fixed.

\paragraph{4. Principal, very small, very large, and off-height terms.}
The new geometry leaves the source contour boxes inside their stated
domains; in particular $\theta>87/100$ and the physical scalar boundary
is still exactly $11/12-\ell/4$. The principal double residue B0.35--B0.36
is the same function $f_\eta$. Its three remainders have relative
reserves $y/4$, $h/24$, and a fixed positive prime-product exponent
$\mu_P$. Their derivation uses the local Euler identity and the
principal correction, not the new moment.

For nonprincipal rows below $Z^{d_0}$ the original global-line bound
B0.56 gives, before small losses,
\[
 h\frac{13}{75}-\frac y2+\frac{63d_0}{50}
 =-\frac{21483987121}{250000000000}<-\frac1{20}.
\]
The finite nonprincipal unit rows are included. For rows above
$Z^{h+\zeta_0}$ the raw absolute estimate B0.57 applies: after choosing
$\zeta_0>0$, choose its fixed positive $z$-line parameter $v$ large
enough that the term has any prescribed negative power. This is not
an application of a central-row density estimate.

The compensated principal probe satisfies the actual B0.27 geometry:
with $M=x+y=1-\ell$ one has
\[
 y-2x+\ell=-\frac{597590039}{5000000000}<0,\qquad
 1+3\ell-2M=-\frac{123299}{750000}<0.
\]
Thus its exponent is $(5x+y)/12=\theta-k$, where
$k=5/6-x/3-\ell/6$. Division by $c_SA(Z)$ costs only a power of
$\log Z$, giving $J_\eta(Z)\ll_\eta Z^{\theta-k+\varepsilon}$.
The three independent contour frequencies and their
joining segments are those of B0.7. Raw polynomial bounds are taken
before any separation in a tail; they give $Z^{B'}T^{-N}$ with $B'$
fixed before the tail order $N$. No contour crosses an unaccounted
nonprincipal numerator pole or a reciprocal zero.

\paragraph{5. Order of choices and continuation.}
Choose from $\mu$, the geometric and certified margins a common moment
loss, capacity decrement, amplitude mesh, and a finite slot system
meeting both the old and new moment meshes on $d\ge d_0$. Its smallest
slot fixes $\mu_P>0$. Choose common real buffers smaller than these
reserves. These choices, including the power saving below, precede
the target $\eta$.

For each target now fix its finite set $S$, dictionary, and the finitely
many internal moment, profile and Sobolev orders. They may depend on the
target. Next choose a target-dependent positive analysis-height exponent
$\tau$ small enough that all factors $(1+Z^\tau)^J$ use less than the
same reserved real margin. The internal detector exponent in the density
proof was already fixed; it is not this later $\tau$. Choose the external
tail order after $\tau$, and finally the target-dependent asymptotic
threshold. The resulting constants and thresholds may depend on the
target, but one obtains a single $\Delta>0$, independent of it, with
\[
 J_\eta(Z)-f_\eta(Z)
       \ll_\eta Z^{\beta_*-k-\Delta}.                \tag{EA.6}
\]
Here $\Delta$ is smaller than fixed fractions of $\mu$, the retained
saving, $y/4,h/24,\mu_P$, and the small-row saving. Make it also smaller
than $\mu/4$. The principal probe bound then gives
$f_\eta(Z)\ll_\eta Z^{\beta_*-k-\Delta}$.

At $Z\downarrow0$ the defining Mellin line can be moved arbitrarily far
right, as in B0.60. Therefore
\[
 \int_0^\infty f_\eta(Z)Z^{-s+k}\,\frac{dZ}{Z}
\]
is holomorphic for $\Re s>\beta_*-\Delta$. On $\Re s>1$ its exact
value is
\[
 e^{(s-5/6)^2}\frac{H_\eta(s)}{L^S(s,\eta)}.
\]
The source cutoff ensures $|H_\eta(s)-1|\le1/2$ for
$\Re s>\theta>87/100$, and the Gaussian never vanishes. Since
$\beta_*-\Delta>\theta$, holomorphic continuation excludes every zero
of $L^S(s,\eta)$ to the right of $\beta_*-\Delta$, away from the
principal pole. Restoring fixed Euler factors does not add zeros in
this positive half-plane.

This is precisely the quantifier pattern of the proved, variable-boundary
\path{Cycle25.PrimitiveProbeContract} and
\path{UniformCommonProbe} in
\path{Cycle25/Assembly/Continuation.lean}. In that file
\path{beta_le_of_common_probe} first chooses a primitive zero
sufficiently near the supremum and then invokes the target's individual
probe. Its Big-O constants and threshold occur after the target
quantifier. Thus no uniform numerical bound on growing target moduli is
needed. A target-dependent power saving would be insufficient; (EA.6)
has explicitly avoided that error. The common positive improvement
contradicts the definition of $\beta_*$. Hence $\beta_*\le\theta$.

\paragraph{6. Primitive reduction and Dirichlet transfer.}
For an imprimitive Hecke character induced by $\eta^*$, the exact
identity T2.1 expresses its $L$-function as $L(s,\eta^*)$ times finitely
many factors $1-\eta^*(\mathfrak p)(N\mathfrak p)^{-s}$.
They are nonzero for $\Re s>0$. This proves the Hecke assertion for
every presented finite-order character of trivial infinite type,
with the same pole exception.

For an ordinary Dirichlet character $\chi$, use its norm lift to the
Eisenstein field. The pinned theorem
\path{HeckeDirichlet.LFunction_eq_dirichlet_product} gives
\[
 L_K(s,\eta_\chi)=L(s,\chi)L(s,\chi\chi_{-3})
\]
away from $s=0,1$, with the original imprimitive and ramified Euler
factors included. For $\theta<\Re s<1$, both factors have finite
values and their product is nonzero by the Hecke result, so
$L(s,\chi)\ne0$. On $\Re s\ge1$ the classical ordinary Dirichlet
nonvanishing theorem applies with its principal-pole exception. The
trivial character of modulus one gives the stated zeta assertion.
These are exactly the transfers in
\path{Cycle25.dirichlet_of_hecke} and
\path{zeta_of_hecke}; they require neither primitivity nor reality
of $\chi$.
\end{proof}

\paragraph{Verification status.}
This is a complete proposed written source-level assembly, under review,
relative to the pinned incumbent analytic framework and the new local
proof. The strict rational comparisons and Bernstein certificate are
exact. Selected supporting Lean modules compile, but this file does not
create a new all-family Lean theorem, a clean dependency replay, a
canonical-target comparison, or an independent kernel certificate.
